\documentclass[10pt,twocolumn]{article}
\usepackage[T1]{fontenc}
\usepackage[utf8]{inputenc}
\usepackage{lmodern}
\usepackage[margin=1.8cm,top=2.4cm,bottom=2cm,columnsep=0.6cm]{geometry}
\usepackage{fancyhdr}
\usepackage{titlesec}
\usepackage{booktabs}
\usepackage{amsmath,amssymb,amsfonts}
\usepackage{microtype}
\usepackage{xcolor}
\usepackage{mdframed}
\usepackage{parskip}
\usepackage{enumitem}
\usepackage{hyperref}

\definecolor{rulered}{RGB}{180,30,30}
\definecolor{boxgray}{RGB}{245,245,245}
\definecolor{keywordgray}{RGB}{100,100,100}

\pagestyle{fancy}
\fancyhf{}
\fancyhead[L]{\textsc{\small Claude Mag}}
\fancyhead[C]{\small\textit{Machine Learning Research Digest}}
\fancyhead[R]{\small 2026-10-26}
\fancyfoot[C]{\small\thepage}
\renewcommand{\headrulewidth}{0.8pt}

\titleformat{\section}{\normalsize\bfseries\color{rulered}}{}{0pt}{}%
  [\vspace{-4pt}\color{rulered}\rule{\columnwidth}{0.4pt}]
\titlespacing{\section}{0pt}{8pt}{4pt}

\hypersetup{colorlinks=true,urlcolor=rulered,linkcolor=black}

\begin{document}
\setlength{\parindent}{0pt}
\setlength{\parskip}{4pt}

{\small\color{keywordgray}\textbf{Keywords:}
  multimodal LLM evaluation \textbullet{}
  visual cognition \textbullet{}
  psychometric benchmark \textbullet{}
  spatial reasoning}\par\vspace{4pt}

{\LARGE\bfseries\raggedright
  Human Cognitive Benchmarks Reveal Foundational Visual Gaps
  in MLLMs}\par\vspace{4pt}

{\large\itshape\raggedright
  VisFactor: 23 Frontier Models Score Below 31\% on Psychometric Visual
  Tests Normed to Human Populations}\par\vspace{6pt}

{\small\textbf{By} Jen-Tse Huang, Dasen Dai, Jen-Yuan Huang, Youliang Yuan,
  Xiaoyuan Liu, Wenxuan Wang, Wenxiang Jiao, Pinjia He, Zhaopeng Tu,
  Haodong Duan (CUHK-ARISE)}\par\vspace{2pt}

{\small\color{keywordgray}arXiv:2502.16435 \textbullet{}
  NeurIPS 2026}
\par\vspace{2pt}\noindent{\color{rulered}\rule{\columnwidth}{0.6pt}}\vspace{6pt}

Multimodal LLMs score impressively on academic visual benchmarks, yet
those benchmarks allow language priors and common-sense shortcuts to
compensate for weak raw visual perception. VisFactor digitizes 20
vision-centric subtests from the Factor-Referenced Cognitive Tests
(FRCT)---a decades-old psychometric battery designed to isolate specific
human visual faculties---and reveals that 23 frontier MLLMs, including
GPT, Gemini, and Qwen families, achieve a best score of 30.17 out of
100. The human norm is approximately 80. Neither scale nor
chain-of-thought prompting closes the gap, exposing foundational
deficits in spatial reasoning and visualization that conventional
benchmarks conceal.

\begin{mdframed}[backgroundcolor=boxgray,linecolor=rulered,linewidth=1pt,
    innerleftmargin=6pt,innerrightmargin=6pt,
    innertopmargin=6pt,innerbottommargin=6pt]
\textbf{\small AT A GLANCE}\par\vspace{4pt}
\begin{description}[leftmargin=0pt,labelsep=4pt,itemsep=2pt,parsep=0pt,
    font=\small\bfseries]
  \item[Problem:] Academic MLLM benchmarks allow language shortcuts and
    miss the low-level visual cognition that human vision relies on.
  \item[Why it matters:] Systems deployed in visually demanding roles
    (medical imaging, robotics, design) require the cognitive visual
    capabilities that current benchmarks do not probe.
  \item[Method:] VisFactor digitizes 20 FRCT vision subtests spanning
    spatial reasoning, perceptual speed, visualization, and visual memory
    into zero-shot multiple-choice questions.
  \item[Result:] Best of 23 frontier MLLMs scores 30.17/100; human
    norm $\approx$ 80; no model reaches the human level on any domain.
  \item[Evidence:] 23 MLLMs (GPT, Gemini, Claude, LLaMA, Qwen, InternVL);
    ablations confirm visual engagement, rule out position bias.
\end{description}
\end{mdframed}

\section{The Big Picture}

The standard narrative in multimodal AI is one of rapid convergence: each
new MLLM generation narrows the gap to human performance on VQA, image
captioning, and visual reasoning benchmarks. This narrative is misleading
because those benchmarks are designed by ML researchers and systematically
admit language-based strategies that bypass genuine visual computation.

Cognitive psychologists have known for decades that human visual cognition
is a multi-factor system encompassing spatial visualization, perceptual
speed, form completion, and short-term visual memory---each measurable
independently with psychometric instruments normed to large human populations.
VisFactor applies those instruments to MLLMs and finds a strikingly different
picture: all frontier models score far below the human norm on every domain,
with the largest gaps in 3D-to-2D spatial transformations and form completion.

\section{Why Existing Benchmarks Fall Short}

\textbf{Language shortcuts:}
Many VQA questions can be answered partially from the question text alone;
models with strong language priors achieve inflated scores without parsing
the image carefully.

\textbf{Coarse granularity:}
Aggregate accuracy on tasks like MMBench or MMMU merges many heterogeneous
subtasks, making it impossible to identify which specific visual process is
failing.

\textbf{Saturation and ceiling effects:}
Frontier models now saturate several prominent benchmarks, leaving no
headroom to distinguish capability levels among top systems.

\textbf{Lack of psychometric grounding:}
Existing benchmarks carry no normative calibration against human populations,
so ``state-of-the-art'' accuracy has no reference point against actual human
visual cognitive ability.

\section{Core Mechanism: VisFactor}

VisFactor digitizes 20 vision-centric subtests from the FRCT battery, grouped
into four domains:

\begin{description}[leftmargin=0pt,labelsep=6pt,itemsep=2pt,parsep=0pt,
    font=\small\bfseries]
  \item[Spatial Reasoning:] Card Rotations, Paper Folding---tests of mental
    rotation and spatial visualization in 3D.
  \item[Perceptual Speed:] Finding A's, Number Comparison---tests of rapid
    figure-ground segregation and symbol matching.
  \item[Visualization:] Surface Development, Form Board---tests of 3D-to-2D
    reasoning and shape completion.
  \item[Visual Memory:] Shape Memory, Position Memory---tests of short-term
    retention of visual configurations.
\end{description}

Each item is digitized as a multiple-choice question. The FRCT scoring key
provides ground-truth answers. No image captioning or text-only reasoning
is possible on any item; answering requires genuine visual computation.

\section{Technical Formulation}

Let $M$ be an MLLM, $\mathcal{Q} = \{(q_i, A_i, a_i^*)\}_{i=1}^N$ the
question set (image $q_i$, options $A_i$, correct answer $a_i^*$). The
VisFactor score is:
\[
\text{Score}(M) = \frac{1}{N}\sum_{i=1}^N \mathbf{1}\!\left[M(q_i, A_i) = a_i^*\right] \times 100
\]
Domain subscores partition the sum over the four cognitive domains. Random
baseline performance is 25 (four options). Human FRCT norm is approximately 80.
The gap $\Delta(M) = 80 - \text{Score}(M)$ measures how far a model is
from average human performance.

\section{Learning or Inference Procedure}

VisFactor is a zero-shot evaluation protocol:

\begin{enumerate}[itemsep=2pt,parsep=0pt]
  \item Present the model with the item image and the four multiple-choice
    options.
  \item Use a fixed instruction: ``Which option is correct? Answer with a
    single letter.''
  \item Record the model's response; compare to the FRCT scoring key.
  \item Compute domain subscores and total score.
\end{enumerate}

A chain-of-thought variant instructs models to explain their reasoning before
answering. Both standard and CoT protocols are reported.

\section{What the Guarantee Says}

VisFactor's psychometric grounding provides several validity properties:

\begin{itemize}[itemsep=2pt,parsep=0pt]
  \item \textbf{Construct validity:} FRCT subtests were validated against
    factor-analytic studies of human visual cognition; each subtest measures
    a theoretically distinct visual factor.
  \item \textbf{Convergent validity:} High within-domain correlations across
    models confirm that subtests in the same domain measure a common latent
    visual ability in MLLMs.
  \item \textbf{No language-shortcut confound:} Removing the image reduces
    all models to chance (25\%), confirming that scores reflect visual
    engagement rather than language priors.
\end{itemize}

\section{Experimental Findings}

\begin{itemize}[itemsep=2pt,parsep=0pt]
  \item \textbf{Best model:} Scores 30.17/100; human norm $\approx 80$.
    All 23 models score between 25 (chance) and 31.
  \item \textbf{Domain asymmetry:} Perceptual speed $\approx 45\%$;
    spatial reasoning $\approx 28\%$; visualization $\approx 20\%$;
    visual memory $\approx 22\%$.
  \item \textbf{Scale does not help:} Larger models within each family
    show no consistent improvement on VisFactor, unlike most academic
    visual benchmarks.
  \item \textbf{Chain-of-thought:} Marginal improvement on spatial
    reasoning; slight degradation on perceptual speed (CoT slows the
    rapid comparison process the subtest is designed to measure).
  \item \textbf{No saturation:} No model approaches the human norm on
    any domain, leaving substantial room for future progress.
\end{itemize}

\section{Ablations and Interpretation}

\textbf{Image removal:}
All models score $\approx 25\%$ when the image is removed and only the
question text and options are provided, confirming genuine visual engagement.

\textbf{Option order shuffling:}
Results are stable under option permutation (variance $< 2$ percentage points),
ruling out position bias as a confound.

\textbf{Item difficulty calibration:}
Items requiring mental rotation of complex 3D figures are the hardest;
items requiring simple symbol comparison are easiest---matching the
psychometric structure designed for humans. The cross-model difficulty
ordering correlates with human item difficulty at $r = 0.73$.

\textbf{Closed-source vs.\ open-source models:}
Closed-source models (GPT, Gemini, Claude) score marginally higher than
open-source models of comparable size, but the gap is small relative
to the distance from the human norm.

\section*{Reference}

Jen-Tse Huang, Dasen Dai, Jen-Yuan Huang, Youliang Yuan, Xiaoyuan Liu,
Wenxuan Wang, Wenxiang Jiao, Pinjia He, Zhaopeng Tu, Haodong Duan.
\textbf{Human Cognitive Benchmarks Reveal Foundational Visual Gaps in MLLMs}.
arXiv:2502.16435. NeurIPS 2026.
\url{https://arxiv.org/abs/2502.16435}

\end{document}
